% Folder 134-2, batch 08 — pages 141–160 (position in the folder; see below).
% Pass: Opus 5.5 (claude-opus-5-5), 2026-10-15 — first pass, unchecked against the pages by a human.
%
% Special rules for folder 134-2 (issues #44 and #57, the latter as decided by the
% user on 2026-10-15) applied: Pursuing Stacks text is not transcribed; his own
% unprinted letters are transcribed in full; texts by others not addressed to him
% are summarised, one \page{} and one \note{} per page, any mark in his hand in full.
%
% Numbering: \page{} numbers the position in the folder, as in batches 6 and 7 (and
% as the tiles and the facsimile do). From p. 115 on the archivists' pencil runs one
% ahead of the position (no sheet pencilled 115 is in the scan): here p. 141 is
% pencilled 142, p. 160 is pencilled 161. Each page's \note{} gives the pencil.
%
% Contents of the batch:
%  - pp. 141-143 (pencilled 142-144): end of R. Brown's typed English translation
%    « Translation of an extract from a Letter from A. Grothendieck to L. Breen
%    17-19 July, 1975 », his pp. 14-16, signed « Alexandre Grothendieck. » on p. 143.
%    Others' text, not addressed to him: summarised, one note per page. No mark in
%    Grothendieck's hand.
%  - p. 144 (pencilled 145): blank; absent.
%  - pp. 145-160 (pencilled 146-161): the HANDWRITTEN FRENCH ORIGINAL of the letter
%    to Lawrence Breen, « Cher Larry », « Villecun 17/19 juillet 1975 », his pages
%    1 to 16 (numbered by him at the head of each page). Transcribed in full, with
%    the full apparatus. Sections in this batch: opening (congratulations, a post at
%    Montpellier), Breen's non-strict n-categories, Dold-Puppe and groups in
%    (1-cat) as crossed modules (L1, L0, d, theta), groups in (n-cat); biduality and
%    homology: A) the six operations, f^!(Lambda) as singular chains, remarks 1)-4);
%    B) Artin-Mazur autoduality of the jacobian; the Serre correspondence of the late
%    50s; ( ) [the « (C) » of the English copy, letter left blank here] complex of
%    generalised jacobians J_{*X/k}, formula (*); D) duality for p-torsion
%    coefficients on curves, Serre duality of unipotent groups, the formula
%    Delta(Rf_!F) = Rf_*(DF[2]), geometric class field theory and Lang's trick.
%    The letter's TEXT BREAKS at the foot of p. 160 (his p. 16) in the middle of a
%    word, inside the bracket « [sur la relation entre le pi_1 "arithmétique" … à
%    noyau un groupe con- »; it continues in batch 9 (p. 161, his p. 17) and runs on
%    into batch 10 (signed on p. 184).
%    The English typed copy (batch 3, pp. 56-60; batch 4) was used only to check the
%    order of the argument and the letter of item « (C) »; no \ill{} or \uncertain{}
%    here was filled from it. Where the French differs from it (e.g. « il semble que
%    tout s'écroule » struck on p. 157 but kept in the English), the note says so.
%
% Letters in this batch: Brown's translation (pp. 141-143, summarised); his own
% letter to L. Breen, 17/19.7.1975, French autograph (pp. 145-160, in full, continued).
\documentclass[11pt,a4paper]{article}
\input{../preamble/grothendieck}

\folder{134-2}
\batch{8}
\pages{141}{160}
\foldertitle{[Chapitre I : Take off (pages 1 à 65) et table des matières provisoire] : tapuscrits annotés (19/02-22/02/1983), note manuscrite (s.d.), copies de lettre (1975, s.d.).}
\dating{1975-[1983]}
\watermark{Édition de démonstration}

\begin{document}

\section{Traduction anglaise (R. Brown) d'un extrait de la lettre à L. Breen des 17--19.7.1975 (fin)}

\page{141}
\note{numérotation : à partir de la p.~115, le crayon des archivistes a un de plus que la position dans le dossier ; cette page est crayonnée 142. P.~« 14 » de la traduction dactylographiée de R. Brown, texte d'un tiers, résumé. Fin de l'alinéa sur le « type d'homotopie fin » : il contient (à la limite) la connaissance de tout, dépend fonctoriellement de $X$ ; pour un schéma de type fini sur un corps algébriquement clos, les théorèmes de finitude cohomologique les plus forts s'énoncent en termes de type d'homotopie fin, dont les constituants élémentaires sont des « polyèdres finis » (voire des variétés compactes à bord), ou plutôt leurs complétés profinis premiers à $p$ au sens d'Artin--Mazur ; programme clair en caractéristique~0, plus mystérieux pour $p > 0$ (variétés qui ne se relèvent pas, même birationnellement). Puis : pas encore de programme algébrique précis, seulement des remarques marginales ; la « linéarisation » d'un type d'homotopie : quelle information cohomologique (coefficients constants ou tordus, produits $H^i(X,M) \times H^j(X,N) \to H^{i+j}(X, M \otimes N)$, opérations cohomologiques) faut-il pour reconstituer le type d'homotopie. Aucune marque de la main de Grothendieck.}

\page{142}
\note{p.~« 15 » de la traduction. Suite de la question de linéarisation : la sous-catégorie de la catégorie dérivée formée des complexes à faisceaux de cohomologie localement constants, avec sa structure triangulée, $\overset{\mathbb{L}}{\otimes}$ et $\mathbb{R}\mathrm{Hom}$, candidat plus satisfaisant ; en allant au fond de la « linéarisation », c'est-à-dire en travaillant non abélien avec les $(n+1)$-catégories de $n$-champs de faisceaux constructibles, on reconstitue le type d'homotopie par son $\infty$-groupoïde fondamental, « comme expliqué dans ma lettre précédente ». De même, pour une stratification $P$ d'un espace topologique modéré $X$, le type d'homotopie plus compliqué attaché à des diagrammes finis devrait correspondre à la $(n+1)$-catégorie des $n$-champs localement constants sur chaque strate (« subordonnés à $P$ »). Aucune marque de sa main.}

\page{143}
\note{p.~« 16 » de la traduction, fin de l'extrait, signée « Alexandre Grothendieck. ». On attend de retrouver le type d'homotopie mixte à partir de la sous-catégorie de la catégorie dérivée des faisceaux abéliens cohomologiquement localement constants sur toutes les strates, avec $\overset{\mathbb{L}}{\otimes}$, $\mathbb{R}\mathrm{Hom}$ et les quatre opérations $\mathbb{R}g_!$, $\mathbb{R}g_*$, $Lg^*$, $g^!$ pour les inclusions de réunions localement fermées de strates ; mais on ne sait pas encore ce qu'est une catégorie triangulée, ni sa version non commutative, une « catégorie homotopique » avec « fibrations » et « cofibrations » (entre crochets : première approche de Quillen au début des années 70). Fin de la « lettre-fleuve » ; question sur la « merveilleuse formule de Bloch--Quillen ». Aucune marque de sa main.}

\section{Lettre à Lawrence Breen, Villecun, 17/19 juillet 1975 (original manuscrit)}

\page{145}
\note{numérotation : cette page est crayonnée 146 ; la page crayonnée 145 (position 144) est blanche. Original manuscrit, en français, de la lettre dont les pp.~56 à 77 (lots~3 et~4) donnent la version anglaise dactylographiée et les pp.~128 à 143 la traduction de R. Brown. Pagination de l'auteur : « 1 » ; lieu et date en haut à droite.}
Villecun 17/19 juillet 1975

Cher Larry,

Tout d'abord félicitations pour la naissance de ta fille, qui heureusement ne s'est pas laissée impressionner \uncertain{par} les importantes \emph{cogitations} mathématiques de son papa -- au risque de les perturber quelque peu.

Au fait, quelle est ta situation en France : es-tu français ou américain, prof. associé ou titulaire ? Si cela pouvait s'arranger, serais-tu intéressé par une position à Montpellier ? Cela te dit-il quelque chose de venir à l'occasion à Montpellier -- ou à Villecun -- pour discuter de tes \uncertain{fantaisies} avec moi -- pour le moment il n'y a guère que moi ici pour ne pas \uncertain{nager} entièrement (seulement partiellement) dans ce que tu fais.
\note{« es-tu » : « es » surcharge un « est ».}

J'ai été content de finir par recevoir un écho à ma \emph{longue} épître -- et même un début d'approche constructive pour \struck{\ill{}} une théorie du type que je prévois. La construction que tu proposes de la notion de $n$-catégorie pas stricte, et du foncteur nerf, a certes le mérite \add{d'exister et} d'être une première approche précise, mais par ailleurs prête \struck{\ill{}} \add{bien sûr} le flanc à des critiques évidentes : d'être très technique, peu intuitive (alors qu'au niveau des 1-cat. et

\page{146}
\note{crayonnée 147 ; p.~« 2 » de la lettre.}
même des 2-cat. tout paraît si clair qu'il n'y a qu'à « suivre le bout de son nez » …) et enfin l'absence d'une définition d'un foncteur allant des ens. (semi) simpliciaux vers les $n$-groupoïdes. Ce foncteur correspond à une intuition géométrique si claire que une théorie qui ne l'inclurait pas me semblerait canularesque. Peut-être en \add{\uncertain{essayant} d'}écri\struck{vant} (comme une sorte de liste de vœux de Noël !) de façon assez complète et explicite les \struck{relations} notions dont on voudrait disposer et les relations (foncteurs, équivalences etc.) qui doivent les lier, arriverait-on finalement à une sorte de description axiomatique assez poussée, \uncertain{qui} soit donne la clef d'une construction ad-hoc explicite, soit \add{qui permette} à la rigueur d'énoncer et démontrer un th. d'existence et d'unicité d'une théorie du type souhaitée.
\note{« \uncertain{essayant} d' » est écrit au-dessus de « écri- », dont la fin est biffée sans être récrite. Le mot après « poussée » est surchargé ; « qui permette » est ajouté au-dessus de « à la rigueur ».}

Par ailleurs, n'ayant pas compris l'idée de Segal dans ta dernière lettre (que j'ai généreusement envoyée à Illusie …), je ne vois pas comment tu définis les $n$-cat. de Picard -- mais peu importe. Pour ce qui est des $n$-cat. de Picard « strictes », tout ce que je leur demande, c'est qu'elles forment finalement une \add{$(n+1)$-}catégorie $(n+1)$-équivalente à celle des complexes de chaînes de longueur~$n$. D'accord ? Je te remercie

\page{147}
\note{crayonnée 148 ; p.~« 3 » de la lettre.}
d'avoir rectifié dans mon esprit une erreur grossière, due à ma grande ignorance de la top. algébrique et de l'homotopie -- j'étais en effet dans le sentiment que les H-espaces satisfaisant des conditions d'associativité et commutativité assez strictes (disons équivalent à des $\Omega^i X$ avec $i$ arbitrairement grand) correspondaient à des groupes topologiques commutatifs (en m'inspirant de quelques analogies …).
Donc entièrement d'accord avec tes observations p.~5.

Par contre, je suis toujours intrigué par la question suivante : n'y a-t-il pas un analogue du tapis de Dold--Puppe pour les groupes (½)simpliciaux (pas néc. commutatifs) et quelle forme prend-il ? À vrai dire, je considère le yoga
\[
\text{ens. ½ simpliciaux} \longleftrightarrow \infty\text{-groupoïdes}
\]
comme étant essentiellement la version « set-theoretic » ultime de Dold--Puppe, qui s'en déduirait en explicitant seulement que les groupes abéliens dans ($\infty$-cat) « sont » les complexes de chaînes dans (Ab). Il faudrait donc déterminer quels sont les groupes de ($\infty$-cat). Je peux te dire ce qu'il en est dans (1-cat) -- ce sera divulg-
\note{« ½ » note \emph{semi-}. Dans « groupoïdes » de la formule, quelques lettres sont surchargées.}

\page{148}
\note{crayonnée 149 ; p.~« 4 » de la lettre.}
gué en long et en large dans le livre de Mme Sinh, je pense, sous le titre du chapitre des « Gr-catégories \emph{strictes} » (i.e. les \uncertain{isomorphismes} d'associativité, d'unité, et d'inverse $x.x^{-1} = 1$ sont des \emph{identités}). On peut expliciter p.ex. commutativement (via le fait que toute Gr-cat. est Gr-équivalente à une Gr-catégorie stricte) le calcul avec les Gr-catégories se réduit à un calcul très algébrique avec des Gr-catégories \emph{strictes}, par une sorte de « calcul des fractions » (à droite ou à gauche, au choix) du type de celui qui aboutit à la construction des catégories dérivées. Quoi qu'il en soit, voici l'explicitation de cette structure (groupes dans (1-cat.)) en termes de théorie des groupes (1-catégorie dans (Gr)) :
\note{« strictes » est souligné deux fois, la seconde fois d'un trait ondulé. La phrase « On peut expliciter … se réduit » est ainsi sur la page.}

La structure est décrite par un quadruplet $(L_1, L_0, d, \theta)$, avec
\[
L_1 \xrightarrow{\;d\;} L_0
\]
un hom. de groupes ordinaires, et
\[
\theta\colon L_0 \longrightarrow \underline{\mathrm{Aut}}_{\mathrm{gr}}(L_1)
\]
une opération de $L_0$ \struck{\ill{}} sur $L_1$, avec les

\page{149}
\note{crayonnée 150 ; p.~« 5 » de la lettre.}
deux axiomes suivants

a) $d$ commute aux op. de $L_0$, quand on le fait opérer sur $L_1$ par $\theta$, sur lui-même par aut. int. :
\[
d(\theta(x_0).x_1) = \mathrm{int}(x_0)\, d(x_1)
\]

b) \[
\theta(d(x_1)) = \mathrm{int}(x_1)
\]

Ces propriétés impliquent que $\operatorname{Im} d$ est invariant dans $L_0$ (donc $\operatorname{Coker} d = \pi_0$ est défini) et $\pi_1 = \operatorname{Ker} d$ est central dans $L_1$, enfin que $L_0$ opère sur $L_1$ en invariant $\pi_1$, et qu'il opère via $\pi_0$. Le principal invariant \add{coh.} de la situation est évidemment l'inv. de Postnikoff--Sinh
\[
\alpha \in H^3(\pi_0, \pi_1) \,.
\]
\note{« Coker » : les lettres du milieu sont surchargées. « qu'il opère via $\pi_0$ » : ainsi sur la page ; le sens attendu est que l'opération induite sur $\pi_1$ se fait via $\pi_0$.}

J'ai rencontré ces animaux -- sans même les chercher -- dans pas mal de situations, que je ne sais pas te passer en revue ici (je viens d'ailleurs de retomber encore dessus à propos de la classification des groupes formels « ordinaires » sur un corps parfait, en termes de groupes algébriques \emph{affines} de groupes formels \emph{commutatifs}, reliés par des Gr-structures strictes de ce type -- pour expliciter le yoga que « le caractère transcendant des
\marginal{[le formalisme sur un topos quelconque (pas seulement dans (Ens))}
\note{la note marginale est écrite verticalement dans la marge gauche, à la hauteur de « je viens d'ailleurs … corps parfait », avec un crochet et une flèche vers le bas ; son point d'insertion n'est pas marqué. Au-dessus, une ligne plus pâle est illisible : \ill{}. Avant « sur un corps parfait », un mot biffé : \struck{\ill{}}.}

\page{150}
\note{crayonnée 151 ; p.~« 6 » de la lettre.}
groupes formels est concentré en les groupes formels commutatifs, découvert semble-t-il par Dieudonné …). La question que je voudrais soulever est la généralisation aux groupes dans ($n$-cat), on s'attend à trouver des complexes de chaînes non commutatifs
\[
L_n \to L_{n-1} \to \dots \to L_1 \to L_0 \to 1
\]
avec des structures supplémentaires sans doute du type $\theta$, mais lesquelles ? Il est entendu que la signification topologique de ces structures, c'est qu'elles expriment exactement les « types d'homotopie tronqués en dim $\leqslant n$ » des groupes topologiques. \add{(… \uncertain{mesurant} les types d'hom. tr. \struck{\ill{}} en dim $\leqslant n+1$ des espaces top. pointés connexes …)} As-tu des candidats à proposer ?
\note{l'insertion entre crochets est écrite entre les lignes, au-dessous de « tronqués en dim $\leqslant n$ » des groupes, et un trait la rattache à la parenthèse ouverte avant « As-tu » ; son premier mot est lu avec doute. Dans « $\leqslant n$ » le signe est surchargé.}

Tes réflexions de bidualité\add{ et d'homologie}, pour formelles qu'elles soient, se rattachent à une foule de développements, dont certains seulement existent dès à présent, et d'autres demanderaient des développements considérables. Voici les réminiscences que m'ouvrent tes noires questions :
\note{« noires » : lecture de la page ; la traduction de R. Brown omet ce passage (p.~131 : « a commentary on biduality and homology is omitted »).}

\page{151}
\note{crayonnée 152 ; p.~« 7 » de la lettre.}
A) Le formalisme des $Rf_!$, $f^!$ (conjoint à $Rf_*$, $Lf^*$, $\overset{\mathbb{L}}{\otimes}$ et $\mathbb{R}\mathrm{Hom}$, « les six opérations ») porte implicitement en son sein la définition de l'homologie et l'identité essentielle entre formalismes homologiques et cohomologiques. \struck{\ill{}} On dispose pour le moment de ce formalisme pour les faisceaux quasi-cohérents sur les schémas -- séminaire Hartshorne -- pour les espaces top. et les faisceaux de coeff. arbitraires dessus -- Verdier, exposé Bourbaki -- et en coh. étale des schémas, pour des coefficients « discrets » \struck{premiers aux car. résiduelles} \add{(de torsion ou « $\ell$-adiques »)} -- SGA~5 --\add{; enfin pour les faisc. coh. sur les espaces an. complexes -- Verdier -- \uncertain{Ruget} --} (le formalisme reste à développer dans le contexte cristallin, et en car.~0 dans le contexte des modules stratifiés \add{à singularités} à la Deligne, avec éventuellement -- sur le corps $\mathbb{C}$ -- des structures de Hodge à la clef, enfin dans le contexte des motifs ; je suis persuadé qu'il est présent un peu partout encore -- partout, peut-être, où il y a un formalisme de nature cohomologique.) Travaillant en coh. étale sur un schéma \add{séparé} de t.f. sur un corps $k$, disons, avec un anneau de coeff. $\Lambda$ de torsion premier à la car.
\note{« premiers aux car. résiduelles » est barré d'un trait et remplacé par « (de torsion ou “$\ell$-adiques”) », écrit au-dessus ; dans « premiers aux », un mot court est biffé. L'insertion « ; enfin pour les faisc. coh. … » est écrite entre les lignes et rattachée par une accolade ; le dernier nom, après « Verdier », est lu avec doute. « \add{à singularités} » est écrit au-dessus de « à la Deligne ». Le mot biffé en tête de « On dispose » est noirci. La phrase continue p.~152.}

\page{152}
\note{crayonnée 153 ; p.~« 8 » de la lettre.}
le \add{complexe de} faisceaux $f^!(\Lambda_e)$ (où $f\colon X \to \operatorname{Spec} k = e$) joue le rôle du \emph{complexe de chaînes singulières sur $X$ à coeff. dans $\Lambda$}, et $Rf_!(f^!\Lambda_e)$ joue le rôle d'une \emph{homologie} $H_*(X/e)$, vis à vis bien sûr de coefficients sur $e$ qui soient des complexes de $\Lambda$-modules. Tu justifieras aisément cette assertion à l'aide des « th. de dualité globale », par un ou deux tours de passe-passe que je t'épargne ici.

\emph{Remarques} 1) Il n'y a pas à \ill{}, ça marche en toutes dimensions.

2) C'est (au moins du pt de vue de ma philosophie) lié au fait que -- disons pour les coefficients \ill{} \add{(sous condition de « constructibilité »)} d'un th. de « bidualité »\textsuperscript{(*)}, du moins sous réserve de résolution des singularités (mais Deligne m'a dit je crois qu'il sait le prouver sans ça) à valeurs dans un « complexe dualisant » $K_e$ \add{(sur $e$)} $\to K_X$ (sur $X$). Si p.ex. $\Lambda$ est « autodual » \add{(ou Gorenstein)} (p.ex. $\Lambda = \mathbb{Z}/n\mathbb{Z}$), $K_e = \Lambda$, donc \add{le complexe dualisant} $K_X = f^!(K_e)$ n'est autre que le « complexe de chaînes singulières à coeff. dans \uncertain{$\Lambda$} » \ill{}
\note{la dernière ligne est coupée au bas du cliché : « singulières à coeff. dans $\Lambda$ » se lit à moitié, la suite non. Dans la remarque 1), le mot après « à » (« tr…gues » ?) n'est pas lu. Dans la remarque 2), le mot qui suit « coefficients » est enfermé dans une parenthèse et rattaché à l'insertion « (sous condition de “constructibilité”) » écrite au-dessus ; il n'est pas lu. L'appel « (*) » est sa main ; la note correspondante n'est pas sur cette page. La flèche entre $K_e$ et $K_X$ est surchargée ; son sens est celui de la page. « le complexe dualisant » est écrit au-dessus de « $K_X = f^!(K_e)$ ».}

\page{153}
\note{crayonnée 154 ; p.~« 9 » de la lettre.}
3) On peut faire les mêmes choses pour des coeff tels que $\mathbb{Z}_\ell$ (Jouanolou, thèse -- non publiée je crains !)

\add{4) Ça marche pour $f\colon X \to S$ prés. finie séparé, si $f$ a des propriétés de « locale trivialité cohomologique » (propr. locales \emph{en haut}) p.ex. $f$ \emph{lisse} ; on trouve $\underline{H}_{*X/S} = Rf_!\, f^!(\Lambda_S)$.}
\note{la remarque 4) est ajoutée en petits caractères entre les lignes, au-dessus de B).}

B) Artin et Mazur ont étudié, dans un esprit voisin du tien, \emph{l'autodualité} de la jacobienne d'une courbe relative $X/S$. Il faut les consulter pour les résultats précis, peut-être cela marche dès que $X/S$ est propre et plat de dim relative~1 -- en tous cas c'est OK sur un anneau de valuation discrète à fibre \emph{générique} lisse. La fibre spéciale peut être très sauvage. (J'ai utilisé \struck{\ill{}} résultat dans SGA~7 pour prouver, dans le cas des jacobiennes, une conjecture de dualité sur les groupes de composantes connexes associés \add{aux} \struck{des} \emph{V.A.} modèles de Néron de V.A. duales l'une de l'autre …)

\add{(ou début des 60 ?)} Vers la fin des années 50, alors que les grands formalismes cohomologiques ($f_!$, $f^!$, coh. étale etc) sortaient à peine des limbes, le cours de Serre sur la théorie de Rosenlicht et Lang des jacobiennes
\note{un signe en tête de l'alinéa « Vers la fin des années 50 » est noirci. « (ou début des 60 ?) » est écrit au-dessus de « 50 » avec un trait de renvoi. Dans « associés aux des V.A. modèles », « aux » surcharge un mot et « des » est biffé : la phrase se lit « associés aux modèles de Néron ». Le mot biffé devant « résultat » est probablement « beau », lecture non retenue.}

\page{154}
\note{crayonnée 155 ; p.~« 10 » de la lettre, chiffre cerclé.}
généralisées\add{ et le corps des classes « géométrique »} \struck{et plus} (voir le livre de Serre), et plus tard sa théorie « géométrique » du corps de classes \emph{local} à coups de groupes proalgébriques (voir son article sur le sujet), m'avait fait réfléchir à des formulations cohomologiques de ces résultats\add{ et d'autres}, qui devaient être notamment de nature \add{« géom. » telle} que les résultats « arithmétiques » sur un corps de base (ou résiduel) $k$ \add{général (fini, p.ex.)} se ramènent immédiatement par descente au cas « géométrique » du corps de base $\bar{k}$. J'ai eu un échange de lettres avec Serre -- je ne sais si j'en retrouverai des copies, mais je me rappelle qu'il s'agissait de projets de théories assez ambitieux\add{, à coups de résidus généralisés, jacobiennes gén. locales etc,} dans au moins trois directions différentes. Mais je ne suis jamais parvenu, malgré de nombreux essais, à accrocher une bonne volonté à un de ces programmes. Voici quelques mots sur eux :
\note{« et le corps des classes “géométrique” » est écrit au-dessus de « généralisées », dont la suite « et plus » est biffée ; « et d'autres » au-dessus de « résultats » ; « “géom.” telle » au-dessus de « nature » ; « général (fini, p.ex.) » au-dessus de « $k$ » ; « à coups de résidus généralisés, jacobiennes gén. locales etc, » entre les lignes, au-dessus de « dans au moins trois directions ». Le $\bar{k}$ : barre au-dessus de « $k$ », lue ainsi.}

( ) Situation $X$ de type fini sur \emph{corps} $k$, construction d'un complexe de jacobiennes généralisées $J_{*X/k}$ (de long. égale à $\dim X$)
\note{la lettre de l'item est laissée en blanc entre les parenthèses ; la copie anglaise dactylographiée (p.~60, lot~3) porte « (C) », ajouté à la main. Un trait vertical en marge gauche longe les deux dernières lignes.}

\page{155}
\note{crayonnée 156 ; p.~« 11 » de la lettre.}
C'était un complexe de pro-groupes algébriques sur $k$, affines à l'exception de $J_0$ si je me rappelle bien, (\struck{\ill{}} $J_0$ avait comme partie abélienne la partie abélienne de $\underline{\mathrm{Alb}}_{X/k}$ la jacobienne \emph{généralisée} ordinaire). Sa construction, inspirée du complexe résiduel, passait par des jacobiennes généralisées (en un sens \add{cohomologique} \struck{convenable}) des localisés $\operatorname{Spec} \underline{O}_{X,x}$ de $X$ en ses différents points. Il est possible qu'en un premier temps, je me limitais au cas $X$ lisse. Le rôle cohomologique de ce complexe était celui d'un complexe d'\emph{homologie}
\marginal{\ill{} \uncertain{$\underline{H}_0(J_*)$} \ill{} de $X$, i.e. on avait un hom.\ $X \to \underline{H}_0$ universel pour les hom.\ de $X$ dans des groupes loc.\ proalgébriques. Pour $X$ connexe, $\underline{H}_0$ est une ext.\ de $\mathbb{Z}$ par un groupe proalg.\ \uncertain{connexe}.}
\note{la note marginale est écrite verticalement dans la marge gauche, en face des lignes « C'était … en un sens », et isolée par un crochet ; sa première ligne, au bord du cliché, n'est lue qu'en partie. C'est l'alinéa « N.B. » de la copie anglaise (p.~61, lot~4). Avant le second $J_0$, un mot court biffé.}

\[
(*) \qquad H^i(X, G_X) \simeq \mathbb{E}\mathrm{xt}^i(J_{*X/k}, G)
\]
mais pour quels coefficients ? Je pense que je devais \add{prendre des groupes alg. comm. $G$ quelconques} \struck{\ill{}, plus \ill{}, à \ill{}}\add{, mais travailler en top. de Zariski (malédiction !),} \struck{les connexes (brrr !), ou plutôt}\add{ \uncertain{même}} dans le cas $G$ discret, je prenais néanmoins des $H^i$ zariskiens -- cela donne des groupes de cohomologie un peu bêtes, évidemment. Je réalisais qu'il faudrait travailler en coh. étale ultérieurement, et que la construction des $J_{*X/k}$ en serait évidemment modifiée. Quant à la sig-
\note{le signe devant la formule, lu « $(*)$ », est surchargé. Après « je devais », une ligne est biffée d'un trait et remplacée, au-dessus et au-dessous, par les deux insertions ; les mots biffés ne sont lus qu'en partie (« … plus … à … »), de même que « les connexes (brrr !), ou plutôt », barré, que remplace « \uncertain{même} » écrit au-dessus. « malédiction » : ainsi, la copie anglaise porte « malédiction ». L'indice de « $J_{*X/k}$ » à la dernière ligne est surchargé d'un « X ».}

\page{156}
\note{crayonnée 157 ; p.~« 12 » de la lettre.}
nification des $\mathbb{E}\mathrm{xt}^i$ (hypercoh.), à un moment où Serre avait développé le formalisme des groupes proalgébriques, on n'avait pas trop peur de les prendre là-dedans -- au sens d'une catégorie dérivée qui à ce moment là n'avait encore jamais été définie et étudiée en forme … (On a fait quand même un peu de chemin depuis !) J'ai bien l'impression, à la suite de ces antiques cogitations, pour heuristiques qu'elles fussent, qu'on doit pouvoir développer à l'heure actuelle une telle théorie des $J_{*X/k}$, en coh. fppf disons, donnant une formule $(*)$ sans limitation sur le dim $i$ de la cohomologie. (NB mais $J_*$ ne s'arrête évidemment plus en $\dim X = n$, mais en dim $2n$. Il est possible néanmoins que les composants $J_i$ soient de dim~0 pour $i > n$).
\note{« le dim $i$ » : ainsi sur la page, pour « le degré $i$ ». « fppf » est écrit en initiales liées.}

Je crois bien que la construction des $J_*$ ne commute pas au chgt de base, mais que dans la cat. dérivée seulement elle y commute.

D) Soit \struck{$X$} $X/k$ schéma lisse (pour simplifier) sur corps $k$, séparé de t.f.,\add{ de dim.\ relative $d$,} et $n$ entier $>0$. Si $n$ est premier à la caract. et si $F$ est un
\note{la lettre $X$ après « Soit » est surchargée ; au-dessus de « lisse » un signe (« \uncertain{sic} » ?) est noirci. « de dim.\ relative $d$, » est écrit au-dessus de « et $n$ entier ».}

\page{157}
\note{crayonnée 158 ; p.~« 13 » de la lettre.}
faisceau de coeff.\ sur $X$ annulé par $n$, la dualité \emph{globale} nous dit que $Rf_!(F)$ et $Rf_*(R\,\underline{\mathrm{Hom}}(F, \mu_n^{\otimes d}))$ sont duals l'un de l'autre à valeurs dans $(\mathbb{Z}/n\mathbb{Z})_k$, p.ex.\ $Rf_!(\mathbb{Z}/n\mathbb{Z})$ et $Rf_*(\mu_n^{\otimes d})$, ou $Rf_!(\mu_n^{\otimes d})$ et $Rf_*(\mathbb{Z}/n\mathbb{Z})$, sont duals l'un de l'autre -- du moins modulo un shift de l'amplitude $2d$ sur les dimensions. (Comme $\mathbb{Z}/n\mathbb{Z}$ est injectif sur lui-même, cela donne d'ailleurs des dualités parfaites
\[
R^if_!(F) \times R^{2d-i}f_*(\mathbb{R}\underline{\mathrm{Hom}}(F, \mu_n^{\otimes d})) \longrightarrow \mathbb{Z}/n\mathbb{Z}) \,.
\]
\note{les « $f_!$ » des lignes 2 et 11 sont surchargés (point d'exclamation repassé sur un astérisque, semble-t-il). Le $\underline{\mathrm{Hom}}$ de la formule est souligné deux fois.}

Si maintenant $n$ n'est plus supposé premier à la caractéristique, \add{p.ex.\ $n$ une puissance de $p = \mathrm{car}\, k > 0$,} \struck{il semble que tout s'écroule,} à commencer qu'on ne sait plus \add{(pour $d > 1$)} par quoi remplacer $\mu_n^{\otimes d}$ … Le miracle assez extraordinaire, c'est que pour $d = 1$ i.e.\ $X$ une courbe lisse, tout continue à marcher parfaitement, pourvu qu'on l'énonce avec soin ! Les premières vérifications se font p.ex.\ avec $F = \mathbb{Z}/p\mathbb{Z}$, $\mu_p$, $\alpha_p$, $X$ complète -- on trouve que c'est OK grâce essentiellement à l'autodualité de la jacobienne. On peut sophistiquer les exemples, en prenant des coefficients $F$ \emph{tordus}, et $X$ non complète -- on se convainc que ça marche toujours ! Simplement il faut faire attention qu'ici les $R^if_*(F)$, $R^if_!(F)$ ont une structure « continue » (ce sont essentiellement des groupes ind-pro-algébriques)
\note{« p.ex.\ $n$ une puissance de $p = \mathrm{car}\, k > 0$, » est écrit au-dessus de la ligne, et « il semble que tout s'écroule, » est barré d'un trait — la phrase reste ainsi sans verbe principal ; la copie anglaise (p.~62) garde « it seems that everything collapses ». « (pour $d > 1$) » au-dessus de « par quoi ».}

\page{158}
\note{crayonnée 159 ; p.~« 14 » de la lettre.}
cela correspond au phénomène bien connu dans la théorie du corps de classes que la structure du $\pi_{1\,\mathrm{ab}}$ de $X$, quand $X$ est non complète, est \emph{continue} -- donc il en est de même du $H^1(X, \mathbb{Z}/p^n\mathbb{Z})$ disons.

À ce propos, je te \emph{signale} que Serre avait proposé (sans jamais l'écrire, je crois) une théorie de dualité des groupes algébriques commutatifs unipotents\add{ connexes,} \struck{\ill{}} pris \emph{mod isogénie radicielle}, dualité à valeurs dans $\mathbb{Q}/\mathbb{Z}$ (ou $\mathbb{Q}_p/\mathbb{Z}_p$). Il trouve que si (sur $k$ alg.\ clos, disons) $G$ est un tel groupe, alors $G' = \mathrm{Ext}^1(G, \mathbb{Q}/\mathbb{Z})$ peut être muni\add{ canoniquement} d'une structure de groupe quasi-algébrique (i.e.\ défini mod isog.\ radicielle), de façon unique sans doute en vérifiant quelques propriétés de fonctorialité, et en exigeant que pour $G = \mathbb{G}_a$ on trouve sur $\mathrm{Ext}^1(\mathbb{G}_a, \mathbb{Q}/\mathbb{Z}) \simeq \mathbb{G}_a$ la structure habituelle. Soit $\Delta G = G' = \underline{\mathrm{Ext}}^1(G, \mathbb{Q}/\mathbb{Z})$. On trouve $G \simeq \Delta\Delta G$ i.e.\ $\Delta$ est une authentique autodualité ! J'appelle $\Delta$ la \emph{dualité de Serre}. Elle passe bien sûr aux ind-pro-groupes sur un corps de base quelconque (par une alg.\ \add{de car.\ $p > 0$.}) D'ailleurs, pour des groupes finis étales, c'est $\underline{\mathrm{Ext}}^0(G, \mathbb{Q}/\mathbb{Z})$ (dualité de Pontrjagin) qui fournit une dualité parfaite. On peut réunir ensemble, dans une catégorie dérivée convenable, dualité de Serre et dualité de Pontrjagin, en prenant
\note{« connexes, » est écrit au-dessus d'un mot biffé et noirci après « unipotents ». « de car.\ $p > 0$. » est écrit au-dessus de « D'ailleurs » et rattaché à « une alg. » ; la parenthèse « (par une alg.\ …) » se lit ainsi, la fin de « alg. » (« \uncertain{clos} » ?) n'étant pas écrite. « Pontrjagin » : la graphie de la page.}

\page{159}
\note{crayonnée 160 ; p.~« 15 » de la lettre.}
$G \longmapsto \Delta G = \mathbb{R}\underline{\mathrm{Hom}}(G, \mathbb{Q}/\mathbb{Z})$ : on l'appellera encore dualité de Serre (« cohomologique »). Ce sera une magnifique autodualité, si on se place dans une catégorie dérivée où les $\underline{H}^i$ des complexes envisagés sont (à des p.n.g.\ ; à la limite près) des extensions de \emph{groupes} étales par des groupes unipotents connexes. Or on rencontre que de tels complexes, en « intégrant » des coeff.\ finis $F$ sur $X$ par $Rf_!$ ou $Rf_*$. Ceci posé, passant à la limite sur $n$ dans la formulation initiale (en encore partant avec $(\mathbb{Q}/\mathbb{Z})_k$ sur $k$, au lieu de $(\mathbb{Z}/n\mathbb{Z})_k$, et formant $f^!(\mathbb{Q}/\mathbb{Z})_k = (\mu_\infty)_X$), la formule de dualité prend la forme
\[
\boxed{\;\Delta(Rf_!(F)) \simeq Rf_*(DF[2])\;}
\]
\marginal{« shift » des dimensions}
où $D$ est la « dualité de Cartier » $R\underline{\mathrm{Hom}}(F, \mu_\infty)$ (ou $R\underline{\mathrm{Hom}}(F, \mathbb{G}_m)$ au choix ?), et $\Delta$ la dualité de Serre : cohomologie à support propre et à support quelconque s'échangent\add{ par dualité,} quand on prend « en haut » la dualité de Cartin, « en bas » la dualité de Serre.
\note{« (à des p.n.g.\ ; à la limite près) » : lecture incertaine de l'abréviation « p.n.g. » ; elle est transcrite telle quelle. « Or on rencontre que de tels complexes » : ainsi sur la page (pour « Or on ne rencontre que »). La formule est encadrée ; « “shift” des dimensions » est écrit à droite, une flèche pointant sur « $[2]$ ». « Cartin » : ainsi la seconde fois, pour « Cartier ». « par dualité, » est écrit au-dessus de « s'échangent ».}

La validité de la formule de dualité ne fait aucun doute -- le travail principal pour
\note{la phrase continue p.~160.}

\page{160}
\note{crayonnée 161 ; p.~« 16 » de la lettre.}
l'établir consiste certainement dans une description rigoureuse des catégories de coefficients avec lesquels on travaille, aussi bien sur $X$ et sur $k$, et des foncteurs $D$ et $\Delta$ dessus. Comme la définition d'une flèche est immédiate, une fois le travail de \emph{montage} de la machine fait, la validité de la formule doit résulter sans difficulté des dévissages habituels, permettant de nous ramener au cas standard $F = \mathbb{Z}/p\mathbb{Z}$, $\mu_p$, $\alpha_p$ sur $X$ lisse complète (\emph{NB} le cas des coeff.\ premiers à la car.\ est déjà connu.)

Explicitant ce que donne la formule de dualité pour $R^1f_*(G_X)$, où $G$ est un groupe\add{ fini} étale sur $k$ (le cas le plus important étant $G = (\mathbb{Z}/p^n\mathbb{Z})_k$), on retrouve la description à la Serre du « corps de classes “géométrique” », en termes d'extensions par $G$ d'une jacobienne généralisée de $X$. Ainsi, la formule de dualité peut être comprise comme une version cohomologique considérablement enrichie du corps de classes géométrique. Quand le corps de base $k$ est fini, pour retrouver le corps de classes sous sa forme classique, on peut utiliser « l'astuce de Lang » [sur la relation entre le $\pi_1$ « arithmétique » d'un groupe algébrique\add{ comm.} lisse connexe $J$ sur $k$ et son $H^0(k, J) = J(k)$, -- le $\pi_1^{\mathrm{ar}}(J)$ classifiant les isogénies\add{ au-dessus de $J$} à noyau un groupe con-
\note{un crochet ouvre en marge gauche l'alinéa « Explicitant … » ; « fini » est écrit au-dessus de « étale ». Après « relation entre », un signe noirci. « comm. » est écrit au-dessus de « lisse », « au-dessus de $J$ » au-dessus de « à noyau ». Le crochet ouvert devant « sur la relation » ne se ferme pas sur cette page. Le texte de la lettre se poursuit p.~161, au lot~9 : il s'arrête ici au milieu de « con[stant] » (lecture attendue, non portée par la page).}

\end{document}
